%% ═══════════════════════════════════════════════════════════════════════════
\section{Lecture 5: Buckingham's Pi Theorem, Blast Waves, and Rocket-Movie Geometry}
\label{sec:lec05}
%% ═══════════════════════════════════════════════════════════════════════════
\date{07/05/2026}

This lecture turns dimensional analysis from a collection of examples into a
general theorem. Buckingham's Pi theorem tells us how many independent
dimensionless variables are needed to describe a physical problem. We then use
that theorem in two observational problems: estimating the energy of a port
explosion from video frames, and testing whether a visible expanding feature in
a missile-interception movie can really be interpreted as a blast wave.

The central lesson is methodological. Dimensional analysis does not only give
quick formulas. It also tells us how many parameters an experiment must vary,
which variables can be irrelevant, and which interpretation of an observation
is internally consistent.

%%───────────────────────────────────────────────────────────────────────────
\subsection{Buckingham's Pi Theorem}
\label{sec:lec05:buckingham}
%%───────────────────────────────────────────────────────────────────────────

\textbf{The Physical Picture:}
Physical laws cannot depend on our arbitrary choice of units. A relation
between measured quantities must therefore look the same if we switch from
meters to centimeters, seconds to hours, or kilograms to grams. Buckingham's
Pi theorem is the systematic way to convert this unit-invariance into a count
of the truly independent, dimensionless variables.

\subsubsection*{General relation}
Suppose a physical problem contains $N$ dimensional variables
$x_1,\ldots,x_N$. A completely general relation among them can be written as
\begin{equation}
    f(x_1,x_2,\ldots,x_N)=0 .
    \label{eq:lec05:general-relation}
\end{equation}
Assume that the variables use $M$ independent base units
$u_1,\ldots,u_M$. Each variable can be decomposed into a numerical value $v_i$
times powers of the units:
\begin{equation}
    x_i = v_i \prod_{j=1}^{M} u_j^{\alpha_{ij}},
    \label{eq:lec05:unit-decomposition}
\end{equation}
where $\alpha_{ij}$ is the power of unit $u_j$ in $x_i$. For example, if $x_i$
is a velocity and the base units are length and time, then $\alpha_{iL}=1$ and
$\alpha_{iT}=-1$.

\subsubsection*{Unit transformations}
Now rescale the base units by
\begin{equation}
    u_j \longrightarrow u_j^{\mathrm{new}}=\lambda_j u_j .
    \label{eq:lec05:unit-rescaling}
\end{equation}
The physical relation must remain true:
\begin{equation}
    f(x_1^{\mathrm{new}},x_2^{\mathrm{new}},\ldots,x_N^{\mathrm{new}})=0 .
    \label{eq:lec05:invariant-relation}
\end{equation}
To extract the consequence of this invariance, differentiate with respect to
one scale factor $\lambda_k$ and then evaluate at $\lambda_j=1$ for all $j$.
By the chain rule,
\begin{align}
    0
    &=
    \left.\pdv{}{\lambda_k}
    f(x_1^{\mathrm{new}},\ldots,x_N^{\mathrm{new}})
    \right|_{\lambda_j=1}
    \nonumber\\
    &=
    \sum_{i=1}^{N}
    \pdv{f}{x_i}
    \left.
    \pdv{x_i^{\mathrm{new}}}{\lambda_k}
    \right|_{\lambda_j=1}.
    \label{eq:lec05:chain-rule}
\end{align}
Since
\begin{equation*}
    x_i^{\mathrm{new}}
    =
    v_i \prod_{j=1}^{M}(\lambda_j u_j)^{\alpha_{ij}},
\end{equation*}
differentiating with respect to $\lambda_k$ gives
\begin{equation*}
    \left.
    \pdv{x_i^{\mathrm{new}}}{\lambda_k}
    \right|_{\lambda_j=1}
    =
    \alpha_{ik} x_i .
\end{equation*}
Therefore each independent unit gives one differential constraint:
\begin{equation}
    \boxed{
    \sum_{i=1}^{N}
    \alpha_{ik} x_i \pdv{f}{x_i}=0,
    \qquad k=1,\ldots,M .
    }
    \label{eq:lec05:unit-invariance-constraint}
\end{equation}
\textit{Physical significance:} Each independent unit-rescaling removes one
independent dimensional direction from the problem.

\subsubsection*{The theorem}
We began with one relation among $N$ variables. Unit invariance adds $M$
independent constraints, provided the $M$ units are genuinely independent in
the variables of the problem. Thus the physical content can be described using
$N-M$ independent dimensionless combinations:
\begin{equation}
    \boxed{
    F(\Pi_1,\Pi_2,\ldots,\Pi_{N-M})=0 .
    }
    \label{eq:lec05:buckingham-theorem}
\end{equation}
\textit{Physical significance:} Instead of exploring an $N$-dimensional
parameter space, the experiment or theory only needs the $N-M$
dimensionless directions.

The $\Pi_i$ are not unique. If $\Pi_1$ and $\Pi_2$ are valid dimensionless
groups, then $\Pi_1\Pi_2$, $\Pi_1/\Pi_2$, or any non-singular redefinition can
also be used. A useful choice puts the quantity we want to solve for in only
one dimensionless group, so that it can be isolated after the dimensionless
relation is written.

%%───────────────────────────────────────────────────────────────────────────
\subsection{Example: Deflection in a Falling Body}
\label{sec:lec05:coriolis-example}
%%───────────────────────────────────────────────────────────────────────────

\textbf{The Physical Picture:}
If a body is dropped from a high tower, Earth's rotation produces a small
Coriolis deflection. We may not know the exact coefficient without solving the
equations of motion, but dimensional analysis already tells us the possible
parameter dependence.

Let $\Delta x$ be the horizontal deflection, $H$ the height, $g$ the
gravitational acceleration, $\Omega$ Earth's angular velocity, $R_\oplus$
Earth's radius, and $m$ the body's mass. The variables use the dimensions
$M,L,T$, so $N=6$ and $M=3$. We expect three independent dimensionless groups.
A convenient choice is
\begin{equation}
    \Pi_1=\frac{\Delta x}{H},\qquad
    \Pi_2=\frac{H}{R_\oplus},\qquad
    \Pi_3=\frac{g}{\Omega^2 R_\oplus}.
    \label{eq:lec05:coriolis-pi-groups}
\end{equation}
The mass cannot enter any nontrivial dimensionless group because no other
variable carries units of mass. Thus the most general dimensionless relation
can be written as
\begin{equation}
    F\qty(\frac{\Delta x}{H},\frac{H}{R_\oplus},
    \frac{g}{\Omega^2R_\oplus})=0 .
    \label{eq:lec05:coriolis-general-relation}
\end{equation}
Equivalently, after isolating $\Delta x$,
\begin{equation}
    \boxed{
    \Delta x =
    H\,\Phi\qty(\frac{H}{R_\oplus},
    \frac{g}{\Omega^2R_\oplus}) .
    }
    \label{eq:lec05:coriolis-dimensional-result}
\end{equation}
\textit{Physical significance:} Dimensional analysis proves immediately that
the dropped body's mass is irrelevant and that only two dimensionless control
parameters remain.

Physics gives one more piece of information. The Coriolis acceleration is
linear in $\Omega$ and in the falling velocity. In the regime $H\ll R_\oplus$,
the dependence must therefore reduce to
\begin{equation}
    \boxed{
    \Delta x \sim C\,\Omega\,\frac{H^{3/2}}{g^{1/2}},
    }
    \label{eq:lec05:coriolis-scaling}
\end{equation}
where $C$ is a dimensionless coefficient, including the latitude dependence.
\textit{Physical significance:} The theorem gives the allowed dimensionless
form; the equation of motion supplies the leading power of the small rotational
effect.

%%───────────────────────────────────────────────────────────────────────────
\subsection{Strong Blast Waves}
\label{sec:lec05:blast-waves}
%%───────────────────────────────────────────────────────────────────────────

\textbf{The Physical Picture:}
A strong explosion releases energy into a small region. At early times the
pressure and energy density inside the blast are much larger than the ambient
air pressure, so the unperturbed pressure can be neglected. The blast front
then expands supersonically, and the question is how its radius $r(t)$ depends
on time, explosion energy, and air density.

\subsubsection*{Relevant variables}
For a strong spherical blast in a uniform medium, take
\begin{itemize}
    \item $r$: radius of the shock front,
    \item $t$: time since the explosion,
    \item $\rho$: ambient density,
    \item $E$: total released energy that drives the blast.
\end{itemize}
Their dimensions are
\begin{equation*}
    [r]=L,\qquad [t]=T,\qquad [\rho]=ML^{-3},\qquad [E]=ML^2T^{-2}.
\end{equation*}
There are $N=4$ variables and $M=3$ independent units, so Buckingham's theorem
says there is one dimensionless group.

Let
\begin{equation*}
    \Pi=r^w t^x E^y \rho^z .
\end{equation*}
Demanding zero net powers of $L,T,M$ gives
\begin{align}
    L:\quad & w+2y-3z=0, \nonumber\\
    T:\quad & x-2y=0, \nonumber\\
    M:\quad & y+z=0 .
    \label{eq:lec05:blast-exponent-equations}
\end{align}
Keeping $w$ free, the solution is
\begin{equation}
    x=-\frac{2w}{5},\qquad y=-\frac{w}{5},\qquad z=\frac{w}{5}.
    \label{eq:lec05:blast-exponent-solution}
\end{equation}
Choosing $w=1$ gives
\begin{equation}
    \Pi = r\,t^{-2/5}E^{-1/5}\rho^{1/5}.
    \label{eq:lec05:blast-pi}
\end{equation}
Therefore the radius must have the Sedov-Taylor scaling
\begin{equation}
    \boxed{
    r(t)=\Pi'\qty(\frac{E}{\rho})^{1/5}t^{2/5}.
    }
    \label{eq:lec05:sedov-radius}
\end{equation}
\textit{Physical significance:} A stronger explosion expands farther, but only
as $E^{1/5}$; the fifth root makes blast-radius estimates relatively
insensitive to the exact energy.

Solving the same relation for the energy gives
\begin{equation}
    \boxed{
    E=\Pi''\,\rho\,\frac{r^5}{t^2}.
    }
    \label{eq:lec05:sedov-energy}
\end{equation}
\textit{Physical significance:} Video measurements of $r$ and $t$ can estimate
the explosion energy, provided the observed front is still a strong shock.

\subsubsection*{Calibration from a known explosion}
Dimensional analysis does not determine the dimensionless coefficient
$\Pi''$. One way to estimate it is to compare with a known explosion. Using
Trinity as an order-of-magnitude calibration, with $E\sim10^{14}\,\si{J}$,
$r\sim130\,\si{m}$, $t\sim0.025\,\si{s}$, and
$\rho\sim1.2\,\si{kg\,m^{-3}}$,
\begin{equation}
    \boxed{
    \Pi'' \simeq E\,r^{-5}t^2\rho^{-1}\sim 1.4 .
    }
    \label{eq:lec05:blast-calibration}
\end{equation}
\textit{Physical significance:} The calibrated coefficient lets us turn the
scaling law into a practical estimator, while still keeping only order-unity
precision.

%%───────────────────────────────────────────────────────────────────────────
\subsection{Bandar Abbas Port Explosion}
\label{sec:lec05:bandar-abbas}
%%───────────────────────────────────────────────────────────────────────────

\textbf{The Physical Picture:}
On April 26, 2025, a large explosion occurred at the Shahid Rajaee port near
Bandar Abbas, Iran. Public videos show a visible blast front. The question is
whether the frames are close enough to the explosion for the strong-shock
formula to estimate the released energy.

\subsubsection*{Camera far from the explosion}
In one video, the visible front travels about $425\,\si{m}$ in $45$ frames at
$30$ frames per second. The average speed is therefore
\begin{equation}
    v_{\mathrm{avg}}
    \simeq
    425\,\si{m}\,
    \frac{30\,\si{s^{-1}}}{45}
    \simeq 300\,\si{m\,s^{-1}} .
    \label{eq:lec05:camera-one-speed}
\end{equation}
This is already comparable to the speed of sound, so the shock is weak. In
that regime its speed is controlled mainly by the acoustic speed rather than
by the original explosion energy. This video therefore gives a poor energy
estimate.

\subsubsection*{Camera closer to the explosion}
A closer camera gives a displacement of about $69\,\si{m}$ in $4$ frames at
$30$ frames per second:
\begin{equation}
    v_{\mathrm{avg}}
    \simeq
    69\,\si{m}\,
    \frac{30\,\si{s^{-1}}}{4}
    \simeq 5.2\times10^2\,\si{m\,s^{-1}} .
    \label{eq:lec05:camera-two-speed}
\end{equation}
This is significantly above the speed of sound, so the strong-shock
approximation is more plausible. Using Eq.~\eqref{eq:lec05:sedov-energy} with
$r\simeq69\,\si{m}$, $t=4/30\,\si{s}$, $\rho\simeq1.2\,\si{kg\,m^{-3}}$, and
$\Pi''\simeq1.4$,
\begin{equation}
    E
    \simeq
    1.4\,(1.2\,\si{kg\,m^{-3}})
    \frac{(69\,\si{m})^5}{(4/30\,\si{s})^2}
    \simeq
    1.5\times10^{11}\,\si{J}.
    \label{eq:lec05:bandar-energy-joule}
\end{equation}
Since $1\,\si{ton}$ of TNT corresponds to roughly $4.2\times10^9\,\si{J}$,
\begin{equation}
    \boxed{
    E \sim 35\,\si{ton\ TNT},
    \qquad
    \text{with a plausible range } \sim 20\text{--}55\,\si{ton\ TNT}.
    }
    \label{eq:lec05:bandar-energy-tnt}
\end{equation}
\textit{Physical significance:} The result is not a precise forensic
measurement; it is an order-of-magnitude energy estimate controlled mainly by
frame timing and distance calibration.

The sensitivity to measurement errors is easy to understand from
$E\propto r^5t^{-2}$. A moderate fractional error in radius is amplified by a
factor of five in the energy, while a one-frame error in time also matters
strongly when only four frames are used.

%%───────────────────────────────────────────────────────────────────────────
\subsection{Rocket-Movie Geometry and a Consistency Check}
\label{sec:lec05:rocket-movie}
%%───────────────────────────────────────────────────────────────────────────

\textbf{The Physical Picture:}
The lecture then considered a movie of a missile interception seen from
Mevasseret Zion shortly after sunset. The visible expanding feature is bright
against the dark sky because it appears to be sunlit. Its angular radius can
be measured frame by frame. The first tempting interpretation is a blast wave,
but dimensional analysis and geometry allow us to test that interpretation.

\subsubsection*{Weak-blast interpretation}
The measured angular radius is approximately linear:
\begin{equation}
    \theta(t)\simeq \theta_0+\dot{\theta}t,
    \qquad
    \dot{\theta}\simeq1.0046\times10^{-2}\,\si{rad\,s^{-1}} .
    \label{eq:lec05:rocket-angular-fit}
\end{equation}
If the feature is already a weak blast wave, its physical expansion speed is
approximately the local sound speed. Taking $c_s\simeq310\,\si{m\,s^{-1}}$,
\begin{equation}
    \boxed{
    D\simeq\frac{c_s}{\dot{\theta}}
    \simeq
    \frac{310}{1.0046\times10^{-2}}\,\si{m}
    \simeq31\,\si{km}.
    }
    \label{eq:lec05:rocket-weak-distance}
\end{equation}
\textit{Physical significance:} If the observed edge moves at the sound speed,
the angular expansion directly gives the distance to the event.

For an altitude angle $\alpha\simeq35^\circ$, this would imply
\begin{equation}
    h\simeq D\sin\alpha \sim 18\,\si{km}
    \qquad
    \text{or, using horizontal geometry, } h\sim 22\,\si{km}.
    \label{eq:lec05:rocket-weak-height}
\end{equation}
This is a low-altitude answer for a ballistic-missile interception, but not
yet a contradiction by itself.

\subsubsection*{Energy upper bound from early linearity}
If the blast wave is already weak in the first measured frame, then the
Sedov-Taylor strong-shock velocity must have fallen to the sound speed before
or around that radius. From Eq.~\eqref{eq:lec05:sedov-radius},
\begin{equation}
    v=\dv{r}{t}=\frac{2r}{5t}.
    \label{eq:lec05:sedov-velocity}
\end{equation}
Setting $v\sim c_s$ gives $t\sim2r/(5c_s)$. Substituting this into
Eq.~\eqref{eq:lec05:sedov-energy} gives
\begin{equation}
    E \lesssim
    \frac{25}{4}\Pi''\rho c_s^2 r_1^3 .
    \label{eq:lec05:rocket-energy-bound}
\end{equation}
For $\theta_1\simeq1\,\si{mrad}$ and $D\simeq31\,\si{km}$,
$r_1\simeq D\theta_1\simeq31\,\si{m}$. With
$\Pi''\simeq1.4$, $\rho\simeq0.05\,\si{kg\,m^{-3}}$, and
$c_s\simeq310\,\si{m\,s^{-1}}$,
\begin{equation}
    \boxed{
    E \lesssim 1.3\times10^9\,\si{J}
    \simeq0.3\,\si{ton\ TNT}.
    }
    \label{eq:lec05:rocket-energy-upper-limit}
\end{equation}
\textit{Physical significance:} Under the weak-blast interpretation, the
observed early linearity gives only an upper bound on the explosive energy.

\subsubsection*{The sunlight test}
The key extra observation is that the feature is still illuminated by the Sun
after local sunset. If the Sun is $\delta$ below the horizon, a point must be
high enough above Earth's shadow to receive direct sunlight. The geometric
condition is
\begin{equation}
    \boxed{
    h \gtrsim R_\oplus\qty(\sec\delta-1).
    }
    \label{eq:lec05:sunlit-height}
\end{equation}
\textit{Physical significance:} Illumination provides a distance and height
constraint that is independent of the blast-wave dynamics.

For the movie, $\delta\simeq7.6^\circ$, so
\begin{equation}
    h\gtrsim
    6371\,\si{km}\qty(\sec 7.6^\circ-1)
    \simeq57\,\si{km}.
    \label{eq:lec05:sunlit-height-number}
\end{equation}
At an altitude angle $\alpha\simeq35^\circ$, the slant distance is therefore
\begin{equation}
    \boxed{
    D\gtrsim\frac{h}{\sin\alpha}
    \simeq
    \frac{57\,\si{km}}{\sin35^\circ}
    \sim100\,\si{km}.
    }
    \label{eq:lec05:sunlit-distance}
\end{equation}
\textit{Physical significance:} The sunlit geometry pushes the event much
farther away than the weak-blast distance estimate.

If the measured angular radius reaches $\theta\simeq3\,\si{mrad}$, then the
physical radius is at least
\begin{equation}
    r\gtrsim D\theta
    \sim
    (100\,\si{km})(3\times10^{-3})
    \sim300\,\si{m}.
    \label{eq:lec05:rocket-physical-radius}
\end{equation}

\subsubsection*{Swept-up mass argument}
At $h\sim57\,\si{km}$, a representative atmospheric density is
$\rho_{\mathrm{air}}\sim3\times10^{-4}\,\si{kg\,m^{-3}}$. The mass of air
inside a sphere of radius $r\sim300\,\si{m}$ is
\begin{equation}
    M_{\mathrm{air}}
    \sim
    \frac{4\pi}{3}\rho_{\mathrm{air}}r^3
    \sim
    \frac{4\pi}{3}(3\times10^{-4})(300)^3\,\si{kg}
    \sim3\times10^4\,\si{kg}.
    \label{eq:lec05:swept-air-mass}
\end{equation}
This is tens of tons of air. A compact vapor cloud or gas shell from the
interception could not expand through that much air at almost constant speed
unless its own mass and energy were enormous.

\begin{equation}
    \boxed{
    \text{The visible feature is probably not a classical blast wave.}
    }
    \label{eq:lec05:not-classical-blast}
\end{equation}
\textit{Physical significance:} The best interpretation is a sparse cloud of
sunlit fragments, dust, smoke, or trail material, whose apparent boundary grows
in projection rather than a dense shell sweeping up all the surrounding air.

%%───────────────────────────────────────────────────────────────────────────
\subsection{Summary of the Method}
\label{sec:lec05:method-summary}
%%───────────────────────────────────────────────────────────────────────────

\textbf{The Physical Picture:}
Dimensional analysis is strongest when it is paired with observational checks.
It can produce a scaling law, but it also tells us when that scaling law is no
longer the right model.

\begin{itemize}
    \item Buckingham's theorem reduces $N$ dimensional variables with $M$
    independent units to $N-M$ dimensionless variables.
    \item For a strong blast wave, the only possible scaling is
    $r\propto(E/\rho)^{1/5}t^{2/5}$.
    \item A video can estimate an explosion energy only if the observed front
    is still a strong shock.
    \item In the Bandar Abbas example, a close camera gave a plausible estimate
    of order tens of tons of TNT.
    \item In the rocket movie, sunlight geometry contradicted the naive
    weak-blast distance estimate, pushing the interpretation toward sparse
    sunlit debris rather than a classical blast wave.
\end{itemize}

